\documentclass[
  spanish,
  letterpaper, oneside
]{article}

\def\documenttitle {Nociones básicas del derecho constitucional}
\def\documentsubtitle {Constitución y derechos humanos}
\def\documentsubject {Actividad 1 - Garantías Constitucionales}

\def\documentauthor {Martin Jonathan de la Cruz}
\def\coursename {Garantías Constitucionales}
\def\coursecode {025}

\def\universityname {Universidad Abierta y a Distancia de México}
\def\universityfaculty {Licenciatura en Derecho}
\def\universitydepartment {Garantías Constitucionales}
\def\universitydepartmentimage {departamentos/UnADM}
\def\universitydepartmentimagecfg {height=1.57cm}
\def\universitylocation {Roma Norte, Ciudad de México}

\def\authortable {
  \begin{tabular}{ll}
    \textbf{Alumno:}
    & \begin{tabular}[t]{l}
      Martin Jonathan de la Cruz \\
    \end{tabular} \\
    Matrícula: & ES2611202040 \\
    Figura docente: & Aarón López Pérez \\
    Semestre/Bloque: & 2 / 1 \\
    & \\
    \multicolumn{2}{l}{Fecha de realización: \today} \\
    \multicolumn{2}{l}{\universitylocation}
  \end{tabular}
}

\input{template}
\usepackage{helvet}
\renewcommand{\familydefault}{\sfdefault}
\usepackage{array}
\usepackage{longtable}
\usepackage{pdflscape}
\usepackage{tikz}
\usetikzlibrary{arrows.meta,positioning,calc,matrix,fit,shapes.geometric,shadows.blur}
\setcitestyle{authoryear,open={(},close={)}}

% Paleta institucional para el bloque del foro (definida localmente por si la
% plantilla no expone estos nombres de color).
\definecolor{unadmGreenDark}{HTML}{174A3A}
\definecolor{unadmGreen}{HTML}{5F8F3A}
\definecolor{unadmGold}{HTML}{B88A2A}
\definecolor{unadmPaper}{HTML}{F6F7F2}
\definecolor{unadmInk}{HTML}{1F2A24}

% -----------------------------------------------------------------------------
% Bloque de participación del foro (producto): caja institucional con
% - texto JUSTIFICADO en fuente normal (no monoespaciada);
% - numeración de intervenciones a la IZQUIERDA (lista enumerada), sin flechas de
%   salto de línea; el texto fluye y se justifica dentro de la caja;
% - contenido SELECCIONABLE y COPIABLE, más un botón de copia (adjunto .txt
%   embebido en el PDF mediante attachfile).
% -----------------------------------------------------------------------------
\usepackage[most]{tcolorbox}
\usepackage{attachfile}
\usepackage{enumitem}

% attachfile: el 'color' debe darse como TRIPLE RGB numérico (0..1), no como un
% nombre \definecolor[HTML] (eso provoca 'Argument of \atfi@textcolor has an extra }').
\attachfilesetup{color={0.373 0.561 0.227},print=false}

% Botón de copia: incrusta un .txt con el contenido íntegro del foro y ofrece un
% enlace clicable para extraer (copiar) el texto desde el visor de PDF.
% El appearance del propio attachfile (icono PushPin) se oculta con appearance={};
% el texto/estilo visible del botón va en el 2.o argumento y se ve siempre.
\newcommand{\foroCopyButton}[1]{%
  \textattachfile[appearance={},description={Participacion del foro (texto copiable)},mimetype=text/plain]{#1}{%
    \colorbox{unadmGreen}{\small\textcolor{white}{\,Copiar participación\,}}%
  }%
}

% Entorno visual del foro
\newtcolorbox{forobox}[1][]{%
  enhanced,
  breakable,
  colback=unadmPaper,
  colframe=unadmGreen,
  boxrule=0pt,
  leftrule=3pt,
  arc=1.2mm,
  left=4mm, right=4mm, top=3mm, bottom=3mm,
  fonttitle=\bfseries\color{white},
  coltitle=white,
  colbacktitle=unadmGreenDark,
  attach boxed title to top left={xshift=4mm,yshift=-2mm},
  boxed title style={colframe=unadmGreenDark,arc=0.6mm},
  #1
}

\def\coverwatermarkenabled {true}
\def\coverwatermarkimage {img/departamentos/UnADM.pdf}
\def\coverwatermarkopacity {0.16}
\def\coverwatermarkwidth {0.72\paperwidth}
\def\coverwatermarkxshift {0cm}
\def\coverwatermarkyshift {-0.35cm}

\newcommand{\insertcoverwatermark}{%
  \ifthenelse{\equal{\coverwatermarkenabled}{true}}{%
    \AddToShipoutPictureBG*{%
      \begin{tikzpicture}[remember picture,overlay]
        \node[
          opacity=\coverwatermarkopacity,
          inner sep=0pt,
          xshift=\coverwatermarkxshift,
          yshift=\coverwatermarkyshift
        ] at (current page.center) {%
          \includegraphics[width=\coverwatermarkwidth]{\coverwatermarkimage}%
        };
      \end{tikzpicture}%
    }%
  }{}%
}

\begin{document}

\insertcoverwatermark
\templatePortrait
\templatePagecfg
\onehalfspacing

\begin{abstractd}
  Las nociones básicas del derecho constitucional articulan la comprensión inicial de las garantías: la Constitución como norma suprema, la distinción entre derecho objetivo y subjetivo, entre derecho sustantivo y adjetivo, la diferencia entre proceso y procedimiento, y la relación entre derecho humano y garantía. Este escrito examina esas categorías, sostiene respuestas propias contrastadas con la Constitución y la Ley de Amparo, y muestra cómo se enlazan para explicar la defensa jurídica de los derechos en el sistema mexicano.
\end{abstractd}

\templateIndex
\templateFinalcfg

\section{Introducción}

El estudio de las garantías constitucionales exige, antes que cualquier análisis avanzado, precisar un conjunto de nociones básicas que suelen confundirse: qué es la Constitución y por qué es suprema, en qué se distinguen el derecho objetivo y el subjetivo, cómo se separan el derecho sustantivo y el adjetivo, qué diferencia hay entre proceso y procedimiento, y cómo se relacionan el derecho humano y la garantía. Estas categorías no son términos aislados: articulan la manera en que el orden jurídico reconoce derechos, limita al poder y ofrece vías para hacerlos exigibles.

Delimitar con claridad estos conceptos permite comprender la defensa jurídica de los derechos humanos, las garantías para su protección y los medios de control constitucional del sistema mexicano. Una definición imprecisa —por ejemplo, tratar la garantía como sinónimo del derecho protegido— compromete después el análisis de figuras como el juicio de amparo.

Para no apoyar las definiciones en una memoria informal, las afirmaciones normativas centrales se contrastan con fuentes formales: la \textit{Constitución Política de los Estados Unidos Mexicanos} y la Ley de Amparo vigente, ambas publicadas por la Cámara de Diputados del H. Congreso de la Unión \citep{cpeum2026,leyAmparo2026}. Sobre esa base conceptual y normativa se organizan las respuestas que siguen.

\section{Constitución, derechos humanos y garantías constitucionales}

\subsection{Supremacía constitucional y artículo 1.\textsuperscript{o}}
La \textbf{Constitución} es el parámetro superior que organiza el poder público, reconoce derechos y limita la actuación de las autoridades. El artículo 1.\textsuperscript{o} constitucional establece que todas las personas gozarán de los derechos humanos reconocidos en la Constitución y en los tratados internacionales, y ordena interpretarlos favoreciendo en todo tiempo la protección más amplia \citep{cpeum2026}. Por ello, la supremacía constitucional no solo ordena jerárquicamente las normas: también orienta la validez y la aplicación de las decisiones públicas.

\subsection{Principios de los derechos humanos}
Los \textbf{principios de los derechos humanos} precisan el alcance de la protección constitucional: universalidad, interdependencia, indivisibilidad y progresividad vinculan la actuación estatal y evitan que los derechos se entiendan como facultades aisladas o de protección eventual \citep{cpeum2026}. En particular, la \textbf{no regresión} exige que las autoridades no adopten medidas que disminuyan injustificadamente el nivel de protección ya alcanzado, pues constituye una consecuencia del deber de progresividad previsto en el artículo 1.\textsuperscript{o} constitucional \citep{cpeum2026}.

En consecuencia, el reconocimiento de los derechos debe analizarse junto con las obligaciones de promover, respetar, proteger y garantizar, así como con el deber de preservar y ampliar, cuando corresponda, su protección efectiva.

\subsection{Derecho humano y garantía constitucional}
La distinción entre \textbf{derecho humano} y \textbf{garantía} no es meramente terminológica: el primero identifica la prerrogativa vinculada con la dignidad de la persona, mientras que la segunda alude a los mecanismos jurídicos que permiten protegerla y hacerla efectiva. Mi postura es que la protección constitucional solo resulta real cuando existe una vía exigible para reclamar la afectación; por eso, el juicio de amparo debe estudiarse como un medio procesal articulado con el derecho protegido, no como un derecho humano en sí mismo \citep{cpeum2026,leyAmparo2026}.

% ==============================================================================
% META (uso interno, no visible en el PDF): el producto solicitado es un FORO
% (técnica de las 100 Técnicas Didácticas de la UnADM, apertura -> participación
% -> coordinación -> cierre). El bloque siguiente reproduce, de forma textual e
% ÍNTEGRA, la intervención publicada (no un resumen). El texto que lo rodea NO
% nombra 'la actividad' ni 'el foro diagnóstico' como narrador externo: gravita
% sobre el tema, preparando la intervención (arriba) e interpretándola (abajo).
% El bloque usa tcolorbox con texto JUSTIFICADO, numeración de intervenciones a
% la izquierda y un botón de copia (adjunto .txt embebido) para extraer el texto.
% ==============================================================================

\begin{forobox}[title={Participación publicada en el foro}]
\hfill\foroCopyButton{foro-participacion-Actividad-1.txt}\\[-0.5em]
\phantomsection\label{foro-participacion}

\textbf{Asunto:} Nociones básicas del derecho constitucional\par
\medskip
Hola a todas y todos: comparto mi participación al foro respondiendo, desde mis conocimientos previos, las cinco preguntas guía sobre nociones básicas de derecho.
\medskip

\begin{enumerate}[leftmargin=1.6em,labelsep=0.6em,itemsep=0.5em,label=\textbf{\arabic*)}]
  \item \textbf{¿Qué es y cuál es la importancia de la Constitución?} La Constitución es la norma fundamental del Estado mexicano. En ella se organiza el poder público, se establecen las competencias de las autoridades, se reconocen derechos de las personas y se fijan límites al ejercicio del poder. Su importancia radica en que sirve como parámetro superior de validez para las demás normas: ninguna ley, acto de autoridad o decisión pública debe contrariar su contenido. Desde la perspectiva de las garantías constitucionales, también aporta la base para defender derechos cuando una autoridad los desconoce o los restringe indebidamente: el artículo 1.\textsuperscript{o} obliga al Estado a promover, respetar, proteger y garantizar los derechos humanos, y el artículo 133 confirma la supremacía constitucional dentro del orden jurídico mexicano.

  \item \textbf{¿Qué entiende por derecho subjetivo y derecho objetivo?} El derecho objetivo es el conjunto de normas jurídicas vigentes que regulan la conducta social, organizan instituciones y establecen consecuencias; es el derecho visto como sistema normativo (constituciones, leyes, reglamentos, tratados y demás reglas aplicables). El derecho subjetivo, en cambio, es la facultad, prerrogativa o posibilidad concreta que una persona tiene para exigir algo, hacer algo, abstenerse de algo o reclamar protección con base en una norma. La diferencia central es que el derecho objetivo es la regla general y el derecho subjetivo es la facultad o pretensión que una persona puede ejercer a partir de esa regla.

  \item \textbf{¿Qué entiende por derecho sustantivo y derecho adjetivo?} El derecho sustantivo regula el contenido de fondo de las relaciones jurídicas: define derechos, obligaciones, prohibiciones, instituciones y consecuencias principales. El derecho adjetivo, también llamado procesal, establece las vías, formas y procedimientos para hacer valer el derecho sustantivo ante una autoridad competente: cómo presentar una demanda, qué etapas debe seguir un juicio, qué recursos existen, qué plazos deben respetarse y cómo se ejecuta una resolución. Así, el sustantivo responde a la pregunta «qué derecho u obligación existe» y el adjetivo responde a «cómo se reclama, protege o hace efectivo ese derecho».

  \item \textbf{¿Conoce la diferencia entre proceso y procedimiento?} Sí. El proceso es la relación o estructura jurídica completa mediante la cual se busca resolver un conflicto o una pretensión ante una autoridad; su finalidad es obtener una decisión fundada que dé solución al caso y respete las garantías de audiencia, defensa, legalidad y debido proceso. El procedimiento es la secuencia concreta de actos, pasos, plazos y formalidades que permiten desarrollar ese proceso. En otras palabras, el proceso es el marco jurídico orientado a resolver una controversia y el procedimiento es el camino ordenado que se sigue para llegar a esa resolución.

  \item \textbf{¿Qué entiende por garantía y por derecho humano?} Un derecho humano es una prerrogativa vinculada con la dignidad de la persona; en el sistema constitucional mexicano se reconoce a toda persona y exige obligaciones de respeto, protección, promoción y garantía por parte de todas las autoridades, conforme al artículo 1.\textsuperscript{o} constitucional. Una garantía es un mecanismo jurídico que protege, asegura o permite hacer efectivo un derecho; por ejemplo, el juicio de amparo es un medio procesal de protección frente a normas generales, actos u omisiones de autoridad que pueden vulnerar derechos humanos, conforme a la Ley de Amparo reglamentaria de los artículos 103 y 107 constitucionales. La diferencia puede sintetizarse así: el derecho humano es el bien o facultad protegida y la garantía es la vía que ayuda a hacerlo exigible. Sin mecanismos eficaces, los derechos pueden quedar en una declaración formal sin protección real.
\end{enumerate}
\medskip

\textbf{Pregunta para el grupo:} ¿en qué caso concreto han visto que un derecho reconocido se quede sin protección por falta de una garantía eficaz?\par
\medskip
Saludos y quedo atento a sus comentarios.
\end{forobox}

Leída en conjunto, la intervención evidencia que las cinco nociones forman una cadena: la supremacía constitucional da validez al sistema normativo (derecho objetivo), del que nacen facultades concretas (derecho subjetivo); el derecho sustantivo fija su contenido y el adjetivo lo hace reclamable mediante procesos y procedimientos; y la garantía es lo que convierte al derecho humano en algo exigible y no en mera declaración \citep{cpeum2026,leyAmparo2026}. La pregunta final dirige la reflexión, precisamente, al eslabón decisivo: la eficacia de la garantía.

% ==============================================================================
% SECCIONES COMENTADAS DE TRAZABILIDAD EDITORIAL
% Estas secciones orientan al motor inteligente y a la memoria editorial, pero no
% forman parte del producto final de foro diagnóstico salvo que la consigna las pida.
% \section{Matriz diagnóstica de conceptos}
%
% \begin{longtable}{>{\bfseries}p{0.26\textwidth}p{0.31\textwidth}p{0.31\textwidth}}
% \toprule
% Eje conceptual & Comprensión esperada & Riesgo de confusión \\
% \midrule
% Constitución & Norma suprema que organiza el poder, reconoce derechos y limita a la autoridad. & Reducirla a un documento histórico sin función jurídica actual. \\
% Derecho objetivo/subjetivo & Distinguir sistema de normas y facultad concreta de una persona. & Creer que derecho objetivo significa ``derecho correcto'' o más justo. \\
% Derecho sustantivo/adjetivo & Separar contenido de derechos y obligaciones de las vías para reclamarlos. & Mezclar el fondo del derecho con el trámite para hacerlo valer. \\
% Proceso/procedimiento & Comprender el proceso como estructura jurídica y el procedimiento como secuencia de actos. & Usarlos como sinónimos absolutos sin reconocer su alcance distinto. \\
% Derecho humano/garantía & Diferenciar prerrogativa inherente o reconocida y mecanismo de protección. & Llamar garantía al derecho mismo o asumir que basta reconocer derechos sin medios de defensa. \\
% \bottomrule
% \end{longtable}
%
% \section{Aplicación a un microcaso}
%
% Supóngase que una persona es detenida sin que se le informe el motivo, se le impide comunicarse con su defensa y se retrasa injustificadamente su presentación ante autoridad competente. En ese caso, los derechos humanos involucrados serían, entre otros, la libertad personal, la seguridad jurídica, el debido proceso y el acceso a una defensa adecuada. La Constitución funcionaría como parámetro de validez para examinar si la actuación de la autoridad respetó límites jurídicos \citep{cpeum2026}.
%
% El derecho sustantivo permitiría identificar los derechos afectados y las obligaciones estatales de respeto y protección. El derecho adjetivo indicaría qué vía procesal puede usarse para reclamar la afectación. El proceso sería el marco institucional mediante el cual se pide protección; el procedimiento serían los pasos concretos, plazos, escritos y resoluciones que conducen a una decisión. En un supuesto de acto de autoridad, el juicio de amparo puede operar como medio formal de protección, siempre que se cumplan sus requisitos de procedencia \citep{leyAmparo2026}.
%
% \section{Valor pedagógico del foro diagnóstico}
%
% El foro permite que la figura docente observe no solo si el estudiante memoriza definiciones, sino si distingue conceptos cercanos y puede aplicarlos a situaciones simples. Para que el diagnóstico sea útil, conviene leer las respuestas con criterios comparables: precisión conceptual, capacidad de contraste, uso de ejemplos, identificación de dudas y coherencia argumentativa.
%
% También es importante que el foro se convierta en decisiones didácticas. Si predominan confusiones entre derecho humano y garantía, el curso debe incluir ejemplos de mecanismos de protección. Si la dificultad aparece entre proceso y procedimiento, conviene trabajar un caso guiado con etapas. Si el grupo muestra respuestas muy heterogéneas, puede ser útil elaborar un glosario comparado y una matriz de errores frecuentes.
% ==============================================================================

\section{Conclusión}

Las nociones básicas del derecho constitucional forman una cadena y no una lista de definiciones sueltas: la Constitución fundamenta y limita el poder; el derecho objetivo constituye el sistema de normas del que el derecho subjetivo extrae facultades concretas; el derecho sustantivo fija los contenidos de fondo y el adjetivo abre las vías para reclamarlos mediante procesos y procedimientos; y la garantía es lo que transforma al derecho humano en algo exigible \citep{cpeum2026,leyAmparo2026}. Vistas así, estas distinciones muestran que la protección constitucional solo se realiza cuando derechos, límites al poder y vías procesales eficaces operan de manera articulada.

  El punto que considero más relevante es que la eficacia de un derecho no depende únicamente de su reconocimiento, sino de la articulación entre las obligaciones de las autoridades previstas en el artículo 1.\textsuperscript{o} constitucional y una vía procesal idónea para reclamar su vulneración. Por ello, sostengo que una garantía debe valorarse por su capacidad real de ofrecer protección, siempre que se cumplan los requisitos de procedencia y trámite previstos para el juicio de amparo; no basta con mencionarlo como mecanismo abstracto de defensa \citep{cpeum2026,leyAmparo2026}. Esta perspectiva orienta el estudio posterior de la asignatura hacia la relación entre reconocimiento normativo, exigibilidad y eficacia práctica.\footnote{Para redactar y organizar este documento se utilizó una herramienta de inteligencia artificial con el propósito de ordenar ideas y revisar la redacción; su uso no sustituyó el análisis propio ni la interpretación jurídica, que corresponden al autor.}


\clearpage
\nocite{unadmSitioWeb,unadmMallaDerecho2024,cpeum2026,leyAmparo2026,normasApaReferencias2020}
\bibliography{garantias-constitucionales}

\end{document}
